\subsection{Linear Independence and Bases}

Throughout this subsection, let $\mathcal{V}$ be a vector space over $\mathbb{F}$.

\begin{definition} \label{def:linear-independence-basis-dimension} \hlt{Linear Independence, Bases, and Finite-Dimensionality}
\begin{enumerate}[label=(\roman*)]
\setlength{\itemsep}{0pt}
\item A nonempty finite list $\beta=(\mathbf{v}_1,\ldots,\mathbf{v}_m)$ in $\mathcal{V}$ is \hlt{linearly independent} iff $(\forall\,(a_1,\ldots,a_m)\in\mathbb{F}^m)\,[\sum\limits_{j=1}^{m}a_j\mathbf{v}_j=\mathbf{0}\Rightarrow a_1=\cdots=a_m=0]$.\\
The empty list is linearly independent. A finite list that is not linearly independent is \hlt{linearly dependent}.
\item A list $\beta$ is a \hlt{basis} of $\mathcal{V}$ iff $\beta$ is linearly independent and $\operatorname{span}(\beta)=\mathcal{V}$.
\item $\mathcal{V}$ is \hlt{finite-dimensional} iff $(\exists\,m\in\N_0)(\exists\,\beta\in\mathcal{V}^m)\,\operatorname{span}(\beta)=\mathcal{V}$. Otherwise $\mathcal{V}$ is \hlt{infinite-dimensional}.
\end{enumerate}
\end{definition}

\begin{proposition} \label{prop:uniquecoefficientsindependentlist} \hlt{Unique Coefficients for Independent Lists}\\
Let $\beta=(\mathbf{v}_1,\ldots,\mathbf{v}_m)$ be linearly independent in $\mathcal{V}$, $\mathbf{a}\defeq(a_1,\ldots,a_m)\in\mathbb{F}^m$, $\mathbf{b}\defeq(b_1,\ldots,b_m)\in\mathbb{F}^m$. Then $\sum\limits_{j=1}^{m}a_j\mathbf{v}_j=\sum\limits_{j=1}^{m}b_j\mathbf{v}_j$ iff $\mathbf{a}=\mathbf{b}$.
\begin{proof}
\begin{enumerate}[label=\roman*.]
\setlength{\itemsep}{0pt}
\item[$(\Rightarrow)$] $\sum\limits_{j=1}^{m}(a_j-b_j)\mathbf{v}_j=\mathbf{0}$ by subtracting the two sums. Hence $(\forall\,j\in\{1,\ldots,m\})\,a_j-b_j=0$ by linear independence of $\beta$, so $\mathbf{a}=\mathbf{b}$.
\item[$(\Leftarrow)$] The two sums are equal by equality of the coefficient tuples.
\end{enumerate}
\hlt{Structural Commentary}: Subtracting two representations converts coefficient comparison into the defining zero-relation test.\\
\hlt{Key Insight}: Independence detects equality coefficient by coefficient.
\end{proof}
\end{proposition}

\begin{proposition} \label{prop:basis-unique-coordinate-representation} \hlt{Unique Coordinate Representation}\\
Let $\beta=(\mathbf{v}_1,\ldots,\mathbf{v}_n)$ be a finite list in $\mathcal{V}$. Then $\beta$ is a basis of $\mathcal{V}$ iff $(\forall\,\mathbf{v}\in\mathcal{V})(\exists!\,(a_1,\ldots,a_n)\in\mathbb{F}^n)\,\mathbf{v}=\sum\limits_{j=1}^{n}a_j\mathbf{v}_j$.
\begin{proof}
\begin{enumerate}[label=\roman*.]
\setlength{\itemsep}{0pt}
\item[$(\Rightarrow)$] Every $\mathbf{v}\in\mathcal{V}$ has a coefficient tuple by $\operatorname{span}(\beta)=\mathcal{V}$, and that tuple is unique by Proposition \ref{prop:uniquecoefficientsindependentlist}.
\item[$(\Leftarrow)$] $\operatorname{span}(\beta)=\mathcal{V}$ by the existence clause. If $\sum\limits_{j=1}^{n}a_j\mathbf{v}_j=\mathbf{0}$, then $\sum\limits_{j=1}^{n}a_j\mathbf{v}_j=\sum\limits_{j=1}^{n}0\mathbf{v}_j$, so $(\forall\,j\in\{1,\ldots,n\})\,a_j=0$ by the uniqueness clause. Thus $\beta$ is linearly independent and hence a basis.
\end{enumerate}
\hlt{Structural Commentary}: Spanning supplies coordinate tuples, while independence separates any two such tuples.\\
\hlt{Key Insight}: A basis is precisely a list with universal and unique coordinates.
\end{proof}
\end{proposition}

\begin{proposition} \label{prop:appendingtodependentlist} \hlt{Appending to a Dependent List}\\
Let $\beta=(\mathbf{v}_1,\ldots,\mathbf{v}_m)$ be linearly dependent in $\mathcal{V}$. Then $(\forall\,\mathbf{v}\in\mathcal{V})\,(\mathbf{v}_1,\ldots,\mathbf{v}_m,\mathbf{v})$ is linearly dependent.
\begin{proof}
Choose $(a_1,\ldots,a_m)\in\mathbb{F}^m$ such that $\sum\limits_{j=1}^{m}a_j\mathbf{v}_j=\mathbf{0}$ and $(\exists\,j\in\{1,\ldots,m\})\,a_j\neq0$ by dependence of $\beta$.\\
For $\mathbf{v}\in\mathcal{V}$, $\sum\limits_{j=1}^{m}a_j\mathbf{v}_j+0\mathbf{v}=\mathbf{0}$ is a nontrivial relation on $(\beta,\mathbf{v})$.\\
\hlt{Structural Commentary}: The original relation is embedded into the longer coefficient tuple by assigning zero to the appended entry.\\
\hlt{Key Insight}: An existing dependence relation survives every append.
\end{proof}
\end{proposition}

\begin{definition} \label{def:omitting-vector-from-list} \hlt{Omitting a Vector from a List}\\
Let $\beta=(\mathbf{v}_1,\ldots,\mathbf{v}_m)$ be a nonempty list in $\mathcal{V}$. For $j\in\{1,\ldots,m\}$, write $\beta_{\widehat{j}}\defeq(\mathbf{v}_1,\ldots,\widehat{\mathbf{v}_j},\ldots,\mathbf{v}_m)$; if $m=1$, then $\beta_{\widehat{1}}=()$.
\end{definition}

\begin{proposition} \label{prop:omittingfromindependentlist} \hlt{Omitting from an Independent List}\\
If $\beta=(\mathbf{v}_1,\ldots,\mathbf{v}_m)$ is a nonempty linearly independent list, then $(\forall\,j\in\{1,\ldots,m\})\ \beta_{\widehat{j}}$ is linearly independent.
\begin{proof}
Fix $j\in\{1,\ldots,m\}$, and let $(a_i)_{i\neq j}$ satisfy $\sum\limits_{i\neq j}a_i\mathbf{v}_i=\mathbf{0}$. Set $a_j\defeq0$. Then $\sum\limits_{i=1}^{m}a_i\mathbf{v}_i=\mathbf{0}$, so $(\forall\,i\in\{1,\ldots,m\})\,a_i=0$ by independence of $\beta$. Hence every coefficient in the original relation is zero.\\
\hlt{Structural Commentary}: Padding a sublist relation with a zero coefficient turns it into a relation on the original list.\\
\hlt{Key Insight}: Deletion cannot introduce a new cancellation.
\end{proof}
\end{proposition}

\begin{proposition} \label{prop:adjoiningandremovingvectors} \hlt{Adjoining and Removing Vectors from Lists}\\
Let $\beta=(\mathbf{v}_1,\ldots,\mathbf{v}_m)$ be a list in $\mathcal{V}$, and write $(\beta,\mathbf{v})\defeq(\mathbf{v}_1,\ldots,\mathbf{v}_m,\mathbf{v})$.
\begin{enumerate}[label=(\roman*)]
\setlength{\itemsep}{0pt}
\item If $\beta$ is linearly independent and $\mathbf{v}\notin\operatorname{span}(\beta)$, then $(\beta,\mathbf{v})$ is linearly independent.
\item If $(a_1,\ldots,a_m)\in\mathbb{F}^m$, $\sum\limits_{i=1}^{m}a_i\mathbf{v}_i=\mathbf{0}$, and $j\in\{1,\ldots,m\}$ satisfies $a_j\neq0$, then $\operatorname{span}(\beta_{\widehat{j}})=\operatorname{span}(\beta)$.\\
Hence $\operatorname{span}(\beta)=\mathcal{V}$ implies $\operatorname{span}(\beta_{\widehat{j}})=\mathcal{V}$.
\end{enumerate}
\begin{proof}
\begin{enumerate}[label=(\roman*)]
\setlength{\itemsep}{0pt}
\item Let $\sum\limits_{i=1}^{m}a_i\mathbf{v}_i+a\mathbf{v}=\mathbf{0}$. If $a\neq0$, then $\mathbf{v}=-a^{-1}\sum\limits_{i=1}^{m}a_i\mathbf{v}_i\in\operatorname{span}(\beta)$, contrary to hypothesis.\\
Hence $a=0$, and then $(\forall\,i\in\{1,\ldots,m\})\,a_i=0$ by independence of $\beta$.
\item $\mathbf{v}_j=-a_j^{-1}\sum\limits_{i\neq j}a_i\mathbf{v}_i\in\operatorname{span}(\beta_{\widehat{j}})$ by $a_j\neq0$ and the stated relation.\\
Thus every vector of $\beta$ lies in $\operatorname{span}(\beta_{\widehat{j}})$, so $\operatorname{span}(\beta)\subseteq\operatorname{span}(\beta_{\widehat{j}})$. The reverse inclusion is immediate because $\beta_{\widehat{j}}$ is a sublist of $\beta$.
\end{enumerate}
\hlt{Structural Commentary}: The first argument isolates the appended coefficient; the second solves for a nonzero-coefficient entry.\\
\hlt{Key Insight}: Append outside the current span, or delete an entry already controlled by the others.
\end{proof}
\end{proposition}

\begin{theorem} \label{thm:basis-equivalent-characterizations} \hlt{Basis Characterizations}\\
Let $\beta=(\mathbf{v}_1,\ldots,\mathbf{v}_n)$ be a finite list in $\mathcal{V}$. Then the following are equivalent:
\begin{enumerate}[label=(\roman*)]
\setlength{\itemsep}{0pt}
\item $\beta$ is a basis of $\mathcal{V}$;
\item $(\forall\,\mathbf{x}\in\mathcal{V})(\exists!\,(a_1,\ldots,a_n)\in\mathbb{F}^n)\,\mathbf{x}=\sum\limits_{j=1}^{n}a_j\mathbf{v}_j$;
\item $\operatorname{span}(\beta)=\mathcal{V}$, and no proper sublist of $\beta$ spans $\mathcal{V}$;
\item $\beta$ is linearly independent, and $(\forall\,\mathbf{x}\in\mathcal{V})\ (\beta,\mathbf{x})$ is linearly dependent.
\end{enumerate}
\begin{proof}
\begin{enumerate}[label=(\roman*)]
\setlength{\itemsep}{0pt}
\item[$(i)$] $\Rightarrow (ii)$ by Proposition \ref{prop:basis-unique-coordinate-representation}.
\item[$(ii)$] $\Rightarrow (iii)$. The equality $\operatorname{span}(\beta)=\mathcal{V}$ holds by the existence clause.\\
Suppose a proper sublist $\gamma$ of $\beta$ spans $\mathcal{V}$, and choose an occurrence $\mathbf{v}_j$ omitted from $\gamma$.\\
A representation of $\mathbf{v}_j$ by $\gamma$ extends to a $\beta$-coefficient tuple whose $j$th entry is $0$, whereas the standard coefficient tuple for $\mathbf{v}_j$ has $j$th entry $1$. This contradicts the uniqueness clause.
\item[$(iii)$] $\Rightarrow (iv)$. If $\beta$ were linearly dependent, choose a nontrivial relation $\sum\limits_{j=1}^{n}a_j\mathbf{v}_j=\mathbf{0}$ and an index $j$ with $a_j\neq0$.\\
Then $\beta_{\widehat{j}}$ would span $\mathcal{V}$ by Proposition \ref{prop:adjoiningandremovingvectors}(ii), contrary to minimality. Hence $\beta$ is linearly independent.\\
For $\mathbf{x}\in\mathcal{V}$, choose $\mathbf{x}=\sum\limits_{j=1}^{n}a_j\mathbf{v}_j$ by $\operatorname{span}(\beta)=\mathcal{V}$. Then $(a_1,\ldots,a_n,-1)$ is a nonzero relation tuple for $(\beta,\mathbf{x})$.
\item[$(iv)$] $\Rightarrow (i)$. If $\operatorname{span}(\beta)\neq\mathcal{V}$, choose $\mathbf{x}\in\mathcal{V}\setminus\operatorname{span}(\beta)$. Then $(\beta,\mathbf{x})$ is linearly independent by Proposition \ref{prop:adjoiningandremovingvectors}(i), contrary to (iv).\\
Thus $\operatorname{span}(\beta)=\mathcal{V}$, and $\beta$ is a basis by the independence clause in (iv).
\end{enumerate}
\hlt{Structural Commentary}: The implication cycle passes through unique coordinates, minimal spanning, and maximal independence.\\
\hlt{Key Insight}: A basis is both an irredundant spanning list and an independence-maximal list.
\end{proof}
\end{theorem}

\begin{lemma} \label{lem:replacement-lemma} \hlt{Replacement Lemma}\\
Let $\mathcal{V}\neq\{\mathbf{0}\}$, let $\beta=(\mathbf{u}_1,\ldots,\mathbf{u}_r)$ span $\mathcal{V}$, and let $\mathbf{0}\neq\mathbf{v}=\sum\limits_{i=1}^r c_i\mathbf{u}_i$. For $j\in\{1,\ldots,r\}$, set $\beta_{j\leftarrow\mathbf{v}}\defeq(\mathbf{u}_1,\ldots,\mathbf{u}_{j-1},\mathbf{v},\mathbf{u}_{j+1},\ldots,\mathbf{u}_r)$.
\begin{enumerate}[label=(\roman*)]
\setlength{\itemsep}{0pt}
\item $(\exists\,j\in\{1,\ldots,r\})\,c_j\neq0$.
\item If $c_j\neq0$, then $\operatorname{span}(\beta_{j\leftarrow\mathbf{v}})=\mathcal{V}$.
\item If $c_j\neq0$ and $\beta$ is a basis, then $\beta_{j\leftarrow\mathbf{v}}$ is a basis.
\item If $r\geq2$, $\beta$ is a basis, $k\in\{1,\ldots,r-1\}$, and $\mathbf{v}\notin\operatorname{span}(\mathbf{u}_1,\ldots,\mathbf{u}_k)$, then $(\exists\,j\in\{k+1,\ldots,r\})\ \beta_{j\leftarrow\mathbf{v}}$ is a basis.
\end{enumerate}
\begin{proof}
\begin{enumerate}[label=(\roman*)]
\setlength{\itemsep}{0pt}
\item If $(\forall\,j\in\{1,\ldots,r\})\,c_j=0$, then $\mathbf{v}=\sum\limits_{i=1}^{r}c_i\mathbf{u}_i=\mathbf{0}$, contrary to hypothesis.
\item Assume $c_j\neq0$. Then $\mathbf{u}_j=c_j^{-1}\left(\mathbf{v}-\sum\limits_{i\neq j}c_i\mathbf{u}_i\right)\in\operatorname{span}(\beta_{j\leftarrow\mathbf{v}})$, so $\operatorname{span}(\beta)\subseteq\operatorname{span}(\beta_{j\leftarrow\mathbf{v}})$.\\
Conversely, $\mathbf{v}\in\operatorname{span}(\beta)$, so $\operatorname{span}(\beta_{j\leftarrow\mathbf{v}})\subseteq\operatorname{span}(\beta)$. Hence $\operatorname{span}(\beta_{j\leftarrow\mathbf{v}})=\mathcal{V}$.
\item Assume $c_j\neq0$ and $\beta$ is a basis. Let $\sum\limits_{i<j}b_i\mathbf{u}_i+a\mathbf{v}+\sum\limits_{i>j}b_i\mathbf{u}_i=\mathbf{0}$. Then $\sum\limits_{i\neq j}(b_i+ac_i)\mathbf{u}_i+ac_j\mathbf{u}_j=\mathbf{0}$ by substituting $\mathbf{v}=\sum\limits_{i=1}^{r}c_i\mathbf{u}_i$.\\
Thus $ac_j=0$ and $(\forall\,i\neq j)\,b_i+ac_i=0$ by independence of $\beta$. Hence $a=0$ by $c_j\neq0$, and then $(\forall\,i\neq j)\,b_i=0$.\\
Therefore $\beta_{j\leftarrow\mathbf{v}}$ is linearly independent; it spans $\mathcal{V}$ by (ii), so it is a basis.
\item If $(\forall\,j\in\{k+1,\ldots,r\})\,c_j=0$, then $\mathbf{v}=\sum\limits_{i=1}^{k}c_i\mathbf{u}_i\in\operatorname{span}(\mathbf{u}_1,\ldots,\mathbf{u}_k)$, contrary to hypothesis.\\
\item Choose $j\in\{k+1,\ldots,r\}$ with $c_j\neq0$. Then $\beta_{j\leftarrow\mathbf{v}}$ is a basis by (iii).
\end{enumerate}
\hlt{Structural Commentary}: Solving for a nonzero-coefficient entry preserves span; comparing coefficients preserves independence.\\
\hlt{Key Insight}: A vector may replace precisely an entry on which its expansion genuinely depends.
\end{proof}
\end{lemma}


\begin{lemma} \label{lem:lineardependencelemma} \hlt{Linear Dependence Lemma}\\
Let $\beta=(\mathbf{v}_1,\ldots,\mathbf{v}_m)$ be linearly dependent with $\mathbf{v}_1\neq\mathbf{0}$. Then $(\exists\,j\in\{2,\ldots,m\})\,[\mathbf{v}_j\in\operatorname{span}(\mathbf{v}_1,\ldots,\mathbf{v}_{j-1})\land\operatorname{span}(\beta_{\widehat{j}})=\operatorname{span}(\beta)]$.
\begin{proof}
\begin{enumerate}[label=(\roman*)]
\setlength{\itemsep}{0pt}
\item \hlt{Terminal coefficient}: Choose a nontrivial relation $\sum\limits_{i=1}^{m}a_i\mathbf{v}_i=\mathbf{0}$, and set $j\defeq\max\{i:a_i\neq0\}$.\\
The equality $j=1$ would imply $\mathbf{v}_1=\mathbf{0}$ by $a_1\mathbf{v}_1=\mathbf{0}$ and $a_1\neq0$, contrary to hypothesis. Hence $j\in\{2,\ldots,m\}$.
\item \hlt{Earlier-span representation}: $(\forall\,i>j)\,a_i=0$ by the choice of $j$, and therefore $\mathbf{v}_j=-a_j^{-1}\sum\limits_{i=1}^{j-1}a_i\mathbf{v}_i\in\operatorname{span}(\mathbf{v}_1,\ldots,\mathbf{v}_{j-1})$.
\item \hlt{Span preservation}: This representation rewrites every linear combination of $\beta$ without $\mathbf{v}_j$, so $\operatorname{span}(\beta)\subseteq\operatorname{span}(\beta_{\widehat{j}})$.\\
The reverse inclusion holds because $\beta_{\widehat{j}}$ is a sublist of $\beta$. Hence $\operatorname{span}(\beta_{\widehat{j}})=\operatorname{span}(\beta)$.
\end{enumerate}
\hlt{Structural Commentary}: The last nonzero relation coefficient identifies a later vector in the earlier span that can be removed.\\
\hlt{Key Insight}: Choose the final active coefficient to preserve both order and span.
\end{proof}
\end{lemma}

\begin{theorem} \label{thm:finitedimlenofvec} \hlt{Independent Lists Are No Longer Than Spanning Lists}\\
If $\mathcal{V}$ is finite-dimensional, $\alpha=(\mathbf{u}_1,\ldots,\mathbf{u}_m)$ is linearly independent, and $\gamma=(\mathbf{w}_1,\ldots,\mathbf{w}_n)$ spans $\mathcal{V}$, then $m\leq n$.
\begin{proof}
\begin{enumerate}[label=(\roman*)]
\setlength{\itemsep}{0pt}
\item \hlt{Contradictory regime}: Suppose $m>n$. Then $n\geq1$: if $n=0$, then $\mathbf{u}_1\in\mathcal{V}=\operatorname{span}(\gamma)=\{\mathbf{0}\}$, contrary to independence of $\alpha$.
\item \hlt{Successive replacement}: Construct spanning lists $\gamma_k$ of length $n$ with designated occurrences of $\mathbf{u}_1,\ldots,\mathbf{u}_k$ for $0\leq k\leq n$.\\
Set $\gamma_0\defeq\gamma$. Suppose $k<n$ and $\gamma_k$ has been constructed. Then $\mathbf{u}_{k+1}\notin\operatorname{span}(\mathbf{u}_1,\ldots,\mathbf{u}_k)$ by independence of $\alpha$.\\
Expand $\mathbf{u}_{k+1}$ in the spanning list $\gamma_k$. Some coefficient at a nondesignated position is nonzero; otherwise the expansion would place $\mathbf{u}_{k+1}$ in $\operatorname{span}(\mathbf{u}_1,\ldots,\mathbf{u}_k)$. The vector $\mathbf{u}_{k+1}$ is nonzero by independence of $\alpha$. Replace the identified position by $\mathbf{u}_{k+1}$.\\
The resulting list $\gamma_{k+1}$ spans $\mathcal{V}$ by Lemma \ref{lem:replacement-lemma}(ii), has length $n$, and retains every earlier designated occurrence.
\item \hlt{Terminal list}: The designated occurrences of $\mathbf{u}_1,\ldots,\mathbf{u}_n$ occupy all $n$ positions of $\gamma_n$. Hence $\operatorname{span}(\mathbf{u}_1,\ldots,\mathbf{u}_n)=\operatorname{span}(\gamma_n)=\mathcal{V}$.
\item \hlt{Contradiction}: Since $m>n$, $\mathbf{u}_{n+1}\in\operatorname{span}(\mathbf{u}_1,\ldots,\mathbf{u}_n)$ by the terminal equality, contrary to independence of $\alpha$.
\end{enumerate}
\hlt{Structural Commentary}: Each replacement reserves one independent vector while keeping the spanning-list length and span fixed.\\
\hlt{Key Insight}: After all spanning positions are reserved, no further independent direction can remain.
\end{proof}
\end{theorem}

\begin{theorem} \label{thm:reducing-spanning-lists} \hlt{Reducing Spanning Lists}\\
If $\mathcal{V}$ is finite-dimensional, then every spanning list in $\mathcal{V}$ contains a basis sublist.
\begin{proof}
Let $\beta=(\mathbf{v}_1,\ldots,\mathbf{v}_n)$ be a spanning list of $\mathcal{V}$. If $\beta$ is independent, it is a basis. Otherwise remove one vector without changing span: omit $\mathbf{v}_1$ if $\mathbf{v}_1=\mathbf{0}$; if $\mathbf{v}_1\neq\mathbf{0}$, remove a suitable $\mathbf{v}_j$ by Lemma \ref{lem:lineardependencelemma}. The resulting sublist still spans $\mathcal{V}$ and has length one less.\\
Repeating at most $n$ times reaches a spanning independent sublist, hence a basis.\\
\hlt{Structural Commentary}: Each deletion shortens the list without changing span, so finite descent reaches an independent sublist.\\
\hlt{Key Insight}: Repeatedly remove a certified redundancy until only a basis remains.
\end{proof}
\end{theorem}

\begin{theorem} \label{thm:extending-independent-lists} \hlt{Extending Independent Lists}\\
If $\mathcal{V}$ is finite-dimensional, then every linearly independent list in $\mathcal{V}$ extends to a basis of $\mathcal{V}$.
\begin{proof}
\begin{enumerate}[label=(\roman*)]
\setlength{\itemsep}{0pt}
\item \hlt{Setup}: Let $\beta=(\mathbf{v}_1,\ldots,\mathbf{v}_m)$ be linearly independent in $\mathcal{V}$, and choose a finite spanning list $(\mathbf{w}_1,\ldots,\mathbf{w}_n)$ of $\mathcal{V}$.
\item \hlt{Construction}: Starting from $\delta_0\defeq\beta$, inspect $\mathbf{w}_1,\ldots,\mathbf{w}_n$ in order\\
 At stage $j$, append $\mathbf{w}_j$ exactly when $\mathbf{w}_j\notin\operatorname{span}(\delta_{j-1})$, and call the resulting list $\delta_j$.\\
Every $\delta_j$ is linearly independent by induction using Proposition \ref{prop:adjoiningandremovingvectors}(i) at each append.
\item \hlt{Verification}: $(\forall\,j\in\{1,\ldots,n\})\ \mathbf{w}_j\in\operatorname{span}(\delta_n)$ because each $\mathbf{w}_j$ was either already spanned when inspected or appended, and subsequent stages only append. Thus $\operatorname{span}(\mathbf{w}_1,\ldots,\mathbf{w}_n)\subseteq\operatorname{span}(\delta_n)$ by Theorem \ref{thm:span-smallest-containing-subspace}(iii).\\
Hence $\operatorname{span}(\delta_n)=\mathcal{V}$, and the independent list $\delta_n$ is a basis extending $\beta$.
\end{enumerate}
\hlt{Structural Commentary}: The finite scan preserves independence while absorbing each chosen spanning generator.\\
\hlt{Key Insight}: Append exactly the generators not already controlled by the current span.
\end{proof}
\end{theorem}

\begin{corollary} \label{cor:finite-dimensional-spaces-and-subspaces-have-bases} \hlt{Bases and Subspaces in Finite Dimension}\\
If $\mathcal{V}$ is finite-dimensional, then:
\begin{enumerate}[label=(\roman*)]
\setlength{\itemsep}{0pt}
\item $\mathcal{V}$ has a basis;
\item $(\forall\,\mathcal{U}\leq\mathcal{V})\ \mathcal{U}$ is finite-dimensional.
\end{enumerate}
\begin{proof}
Choose a finite spanning list $\gamma=(\mathbf{w}_1,\ldots,\mathbf{w}_n)$ of $\mathcal{V}$.
\begin{enumerate}[label=(\roman*)]
\setlength{\itemsep}{0pt}
\item A basis of $\mathcal{V}$ is contained in $\gamma$ by Theorem \ref{thm:reducing-spanning-lists}.
\item Fix $\mathcal{U}\leq\mathcal{V}$, and set $\beta_0\defeq()$. Whenever $\operatorname{span}(\beta_k)\neq\mathcal{U}$, choose $\mathbf{u}_{k+1}\in\mathcal{U}\setminus\operatorname{span}(\beta_k)$ and set $\beta_{k+1}\defeq(\beta_k,\mathbf{u}_{k+1})$.\\
Every constructed list is linearly independent by induction using Proposition \ref{prop:adjoiningandremovingvectors}(i).\\
The construction cannot reach $\beta_{n+1}$ by Theorem \ref{thm:finitedimlenofvec}, because $\gamma$ is a spanning list of length $n$.\\Hence it stops at some $k\leq n$ with $\operatorname{span}(\beta_k)=\mathcal{U}$. Thus $\beta_k$ is a finite basis of $\mathcal{U}$, so $\mathcal{U}$ is finite-dimensional.
\end{enumerate}
\hlt{Structural Commentary}: Basis reduction handles $\mathcal{V}$; the ambient length bound forces every subspace construction to terminate.\\
\hlt{Key Insight}: Ambient finite spanning length bounds the size of every independent subspace list.
\end{proof}
\end{corollary}

\begin{definition} \label{def:hamel-basis} \hlt{Independent Sets and Hamel Bases}\\
A subset $I\subseteq\mathcal{V}$ is \hlt{linearly independent} iff every finite list of distinct vectors in $I$ is linearly independent.\\
A subset $\mathcal{B}\subseteq\mathcal{V}$ is a \hlt{Hamel basis} of $\mathcal{V}$ iff $\mathcal{B}$ is linearly independent and $\operatorname{span}(\mathcal{B})=\mathcal{V}$.\\
For $a:\mathcal{B}\to\mathbb{F}$, define $\operatorname{supp}_{\mathcal{B}}(a)\defeq\{\mathbf{b}\in\mathcal{B}:a(\mathbf{b})\neq0\}$. A \hlt{finite $\mathcal{B}$-expansion} of $\mathbf{v}\in\mathcal{V}$ is a function $a:\mathcal{B}\to\mathbb{F}$ with finite support such that $\mathbf{v}=\sum\limits_{\mathbf{b}\in\operatorname{supp}_{\mathcal{B}}(a)}a(\mathbf{b})\mathbf{b}$. For a finite Hamel basis, any ordering of $\mathcal{B}$ is an ordered basis.
\end{definition}

\begin{proposition} \label{prop:hamel-coordinate-representation} \hlt{Hamel Coordinate Representation}\\
If $\mathcal{B}$ is a Hamel basis of $\mathcal{V}$, then every $\mathbf{v}\in\mathcal{V}$ has a unique finite $\mathcal{B}$-expansion.
\begin{proof}
Fix $\mathbf{v}\in\mathcal{V}$. Choose $m\in\N_0$, $(c_1,\ldots,c_m)\in\mathbb{F}^m$, and $(\mathbf{b}_1,\ldots,\mathbf{b}_m)\in\mathcal{B}^m$ such that $\mathbf{v}=\sum\limits_{j=1}^{m}c_j\mathbf{b}_j$ by $\operatorname{span}(\mathcal{B})=\mathcal{V}$.\\
For $\mathbf{b}\in\mathcal{B}$, set $J_{\mathbf{b}}\defeq\{j\in\{1,\ldots,m\}:\mathbf{b}_j=\mathbf{b}\}$, and define $a:\mathcal{B}\to\mathbb{F}$ by $a(\mathbf{b})\defeq\sum\limits_{j\in J_{\mathbf{b}}}c_j$.\\
Then $\operatorname{supp}_{\mathcal{B}}(a)\subseteq\{\mathbf{b}_1,\ldots,\mathbf{b}_m\}$, and $\mathbf{v}=\sum\limits_{\mathbf{b}\in\operatorname{supp}_{\mathcal{B}}(a)}a(\mathbf{b})\mathbf{b}$ by regrouping equal entries. Hence $a$ is a finite $\mathcal{B}$-expansion of $\mathbf{v}$.\\
For uniqueness, let $a,a':\mathcal{B}\to\mathbb{F}$ be finite $\mathcal{B}$-expansions of $\mathbf{v}$, and set $F\defeq\operatorname{supp}_{\mathcal{B}}(a)\cup\operatorname{supp}_{\mathcal{B}}(a')$.\\
Then $F$ is finite and $\sum\limits_{\mathbf{b}\in F}(a(\mathbf{b})-a'(\mathbf{b}))\mathbf{b}=\mathbf{0}$. Thus $(\forall\,\mathbf{b}\in F)\,a(\mathbf{b})=a'(\mathbf{b})$ by independence of $\mathcal{B}$. Both vanish on $\mathcal{B}\setminus F$, so $a=a'$.\\
\hlt{Structural Commentary}: Existence consolidates repeated generators; uniqueness compares their combined finite support.\\
\hlt{Key Insight}: Hamel coordinates are ordinary independent-list coefficients indexed by finite support.
\end{proof}
\end{proposition}

\begin{theorem} \label{thm:hamel-basis-extension} \hlt{Hamel Basis Extension}\\
Assume Zorn's lemma. Then:
\begin{enumerate}[label=(\roman*)]
\setlength{\itemsep}{0pt}
\item if $I\subseteq S\subseteq\mathcal{V}$, $I$ is linearly independent, and $\operatorname{span}(S)=\mathcal{V}$, then $(\exists\,\mathcal{B})\,[I\subseteq\mathcal{B}\subseteq S\land\mathcal{B}\text{ is a Hamel basis of }\mathcal{V}]$;
\item $\mathcal{V}$ has a Hamel basis;
\item every linearly independent subset of $\mathcal{V}$ extends to a Hamel basis of $\mathcal{V}$;
\item every spanning subset of $\mathcal{V}$ contains a Hamel basis of $\mathcal{V}$.
\end{enumerate}
\begin{proof}
\begin{enumerate}[label=(\roman*)]
\setlength{\itemsep}{0pt}
\item Let $\mathscr{P}\defeq\{E:I\subseteq E\subseteq S,\ E\text{ is linearly independent}\}$, ordered by inclusion.\\
Then $I\in\mathscr{P}$. Let $\mathcal{C}\subseteq\mathscr{P}$ be a nonempty chain, and set $U\defeq\bigcup\limits_{E\in\mathcal{C}}E$. Then $I\subseteq U\subseteq S$.\\
Let $m\in\N$, and consider a relation $\sum\limits_{k=1}^{m}a_k\mathbf{u}_k=\mathbf{0}$ among distinct vectors of $U$.\\
Choose $E_1,\ldots,E_m\in\mathcal{C}$ such that $(\forall\,k\in\{1,\ldots,m\})\,\mathbf{u}_k\in E_k$.\\
Some $E_\ast\in\mathcal{C}$ contains $E_1\cup\cdots\cup E_m$ by finite induction on the chain order.\\
Thus $(\forall\,k\in\{1,\ldots,m\})\,a_k=0$ by independence of $E_\ast$, so $U\in\mathscr{P}$.\\
Hence every nonempty chain has an upper bound in $\mathscr{P}$; the empty chain is bounded by $I$.\\
Choose a maximal $\mathcal{B}\in\mathscr{P}$ by Zorn's lemma. Suppose $s\in S\setminus\operatorname{span}(\mathcal{B})$.\\
In a finite relation among distinct vectors of $\mathcal{B}\cup\{s\}$, a nonzero coefficient of $s$ would imply $s\in\operatorname{span}(\mathcal{B})$, while a zero coefficient leaves a relation in $\mathcal{B}$. Hence $\mathcal{B}\cup\{s\}$ is linearly independent, contrary to maximality of $\mathcal{B}$.\\
Therefore $S\subseteq\operatorname{span}(\mathcal{B})$, so $\operatorname{span}(\mathcal{B})=\mathcal{V}$ by $\operatorname{span}(S)=\mathcal{V}$. Thus $I\subseteq\mathcal{B}\subseteq S$, and $\mathcal{B}$ is a Hamel basis.
\item $\mathcal{V}$ has a Hamel basis by (i) with the linearly independent set $I=\emptyset$ and the spanning set $S=\mathcal{V}$.
\item Every prescribed linearly independent subset extends to a Hamel basis by (i) with $S=\mathcal{V}$.
\item Every prescribed spanning subset contains a Hamel basis by (i) with $I=\emptyset$.
\end{enumerate}
\hlt{Structural Commentary}: Chain unions preserve finite independence, and maximality rules out every unspanned element of $S$.\\
\hlt{Key Insight}: Finitary independence is stable under chain unions, so Zorn maximality forces a Hamel basis.
\end{proof}
\end{theorem}
